% =====================================================================
\chapter{Derivadas parciales, gradiente y regla de la cadena}
\label{tema:4}
% =====================================================================

La derivada de una función de una variable mide cuánto cambia la función por unidad de cambio de la variable. Con varias variables hay que elegir \emph{en qué dirección} se mueve uno. Las derivadas parciales miden el cambio en las direcciones de los ejes; el gradiente las reúne en un vector que apunta hacia donde la función crece más deprisa y que es perpendicular a las curvas de nivel. Este es el objeto central del entrenamiento de redes neuronales. La presentación sigue las sesiones~26--27 y~32--38 de 18.02SC~\cite{mit1802}, las §4.3--4.6 de OpenStax~\cite{openstax3}, los capítulos~13 de Strang~\cite{strangcalc} y~8 de Apostol~\cite{apostol2}.

\section{Derivadas parciales}

\begin{definicion}[Derivadas parciales]\label{def:parcial}
Sea $f$ una función de $(x,y)$ definida cerca de $(a,b)$. Las \emph{derivadas parciales} de $f$ en $(a,b)$ son
\begin{align*}
f_x(a,b)&=\pd fx(a,b)=\lim_{h\to0}\frac{f(a+h,b)-f(a,b)}{h},\\
f_y(a,b)&=\pd fy(a,b)=\lim_{h\to0}\frac{f(a,b+h)-f(a,b)}{h},
\end{align*}
cuando existen. Para tres o más variables la definición es análoga.
\end{definicion}

Para calcular $f_x$ se deriva respecto de $x$ tratando $y$ como una constante, con todas las reglas de una variable. Geométricamente, $f_x(a,b)$ es la pendiente de la curva que se obtiene al cortar la gráfica con el plano vertical $y=b$.

\addtocounter{definicion}{1}%

La igualdad de las derivadas cruzadas del ejemplo no es casual.

\begin{teorema}[Schwarz]\label{teo:schwarz}
Si $f_{xy}$ y $f_{yx}$ existen y son continuas en un entorno de $(a,b)$, entonces $f_{xy}(a,b)=f_{yx}(a,b)$.
\end{teorema}

La demostración, que usa dos veces el teorema del valor medio, puede consultarse en Apostol~\cite[cap.~8]{apostol2}. Todas las funciones de este módulo construidas con polinomios, exponenciales, logaritmos y funciones trigonométricas cumplen la hipótesis en el interior de su dominio.

\section{Diferenciabilidad y plano tangente}

La mera existencia de las derivadas parciales es una condición débil: solo informa de lo que ocurre sobre dos rectas. La noción adecuada exige que $f$ se parezca a una función afín en \emph{todas} las direcciones.

\begin{definicion}[Diferenciabilidad]\label{def:diferenciable}
Una función $f$ definida cerca de $\va\in\R^n$ es \emph{diferenciable} en $\va$ si existe un vector $\mathbf g\in\R^n$ tal que
\[
\lim_{\mathbf h\to\vcero}\frac{f(\va+\mathbf h)-f(\va)-\mathbf g\cdot\mathbf h}{\norm{\mathbf h}}=0 .
\]
Es decir, $f(\va+\mathbf h)=f(\va)+\mathbf g\cdot\mathbf h+\varepsilon(\mathbf h)\norm{\mathbf h}$ con $\varepsilon(\mathbf h)\to0$.
\end{definicion}

\begin{teorema}\label{teo:diferenciable}
Si $f$ es diferenciable en $\va$, entonces:
\begin{enumerate}
  \item $f$ es continua en $\va$;
  \item existen las derivadas parciales en $\va$ y $\mathbf g=\nabla f(\va)$, donde el \emph{gradiente} es el vector $\nabla f(\va)=\bigl(f_{x_1}(\va),\dots,f_{x_n}(\va)\bigr)$.
\end{enumerate}
\end{teorema}

El recíproco de (ii) es falso (\cref{eje:4-no-diferenciable}). Una condición suficiente práctica es la siguiente, cuya demostración puede verse en Apostol~\cite[cap.~8]{apostol2}: \emph{si las derivadas parciales de $f$ existen y son continuas en un entorno de $\va$, entonces $f$ es diferenciable en $\va$.}

\begin{definicion}[Plano tangente y aproximación lineal]\label{def:tangente}
Si $f$ es diferenciable en $(a,b)$, el \emph{plano tangente} a la gráfica de $f$ en $(a,b,f(a,b))$ es
\[
z=f(a,b)+f_x(a,b)(x-a)+f_y(a,b)(y-b),
\]
y la \emph{aproximación lineal} es $f(a+\Delta x,b+\Delta y)\approx f(a,b)+f_x(a,b)\Delta x+f_y(a,b)\Delta y$. Con notación de diferenciales, $df=f_x\,dx+f_y\,dy$.
\end{definicion}

\addtocounter{definicion}{1}%

\section{Regla de la cadena}

\begin{teorema}[Regla de la cadena]\label{teo:cadena}
Sea $\vr(t)=(x(t),y(t))$ derivable en $t_0$ y $f$ diferenciable en $\vr(t_0)$. Entonces $w(t)=f(\vr(t))$ es derivable en $t_0$ y
\[
\frac{dw}{dt}=f_x\frac{dx}{dt}+f_y\frac{dy}{dt}=\nabla f(\vr(t_0))\cdot\vr\,'(t_0).
\]
\end{teorema}

Si $x$ e $y$ dependen de dos variables $(u,v)$, se aplica el teorema fijando una de ellas:
\begin{equation}\label{eq:t4-cadena2}
\pd wu=\pd wx\pd xu+\pd wy\pd yu,\qquad \pd wv=\pd wx\pd xv+\pd wy\pd yv .
\end{equation}
La regla se recuerda con un \emph{diagrama de árbol}: cada camino de $w$ a la variable final aporta el producto de las derivadas de sus tramos, y se suman las aportaciones de todos los caminos.

\addtocounter{figure}{1}%

\addtocounter{definicion}{1}%

\section{Gradiente y derivada direccional}

\begin{definicion}[Derivada direccional]\label{def:direccional}
Sea $\vu$ un vector \emph{unitario}. La \emph{derivada direccional} de $f$ en $\va$ según $\vu$ es
\[
D_{\vu}f(\va)=\lim_{t\to0}\frac{f(\va+t\vu)-f(\va)}{t},
\]
la tasa de cambio de $f$ por unidad de longitud recorrida en la dirección de $\vu$. Para $\vu=\vi$ y $\vu=\vj$ se obtienen $f_x$ y $f_y$.
\end{definicion}

\begin{teorema}[Cálculo de la derivada direccional]\label{teo:direccional}
Si $f$ es diferenciable en $\va$, entonces $D_{\vu}f(\va)=\nabla f(\va)\cdot\vu$ para todo vector unitario $\vu$.
\end{teorema}

\begin{corolario}[Dirección de máximo crecimiento]\label{cor:maximo-crecimiento}
Sea $\nabla f(\va)\neq\vcero$. Entre todas las direcciones unitarias, $D_{\vu}f(\va)$ alcanza su máximo, $\norm{\nabla f(\va)}$, para $\vu=\nabla f(\va)/\norm{\nabla f(\va)}$, su mínimo, $-\norm{\nabla f(\va)}$, en la dirección opuesta, y se anula en las direcciones perpendiculares al gradiente.
\end{corolario}

\begin{teorema}[El gradiente es normal a los conjuntos de nivel]\label{teo:gradiente-normal}
Sea $f$ diferenciable y $\vr(t)$ una curva derivable contenida en el conjunto de nivel $f=c$, con $\vr(t_0)=\va$. Entonces $\nabla f(\va)\cdot\vr\,'(t_0)=0$.
\end{teorema}

En consecuencia, el \emph{plano tangente} a la superficie de nivel $F(x,y,z)=c$ en un punto $P_0$ con $\nabla F(P_0)\neq\vcero$ es el plano de normal $\nabla F(P_0)$: $\nabla F(P_0)\cdot(X-P_0)=0$. El plano tangente a una gráfica $z=f(x,y)$ es el caso particular $F(x,y,z)=f(x,y)-z$.

\addtocounter{definicion}{1}%

\addtocounter{definicion}{1}%

\addtocounter{definicion}{1}%

